\scanpage{22}
\mspage{6}
\textcircled{5.}
\uline{Germes de voisinages côniques, ou cylindriques.}
\deleted{\uline{Compatibilité.}}\ \uline{Transversalité.}

Soit $Y\subset X$. Deux voisinages coniques
\inserted{$\bar T,\bar T'$} de $Y$ dans
$X$ sont équivalents s'il existe $\lambda,\lambda'\in]0,1[$ tels
que $\lambda\bar T=\lambda'\bar T'$ (égalité \inserted{compatible}
avec les structures
coniques sur $Y$). C'est une relation d'équivalence.
\margin{$\varepsilon\in]0,1[$, alors pour
$\varepsilon\lambda\bar T=\varepsilon\lambda'\bar T'$}
Une classe d'équivalence est un \uline{germe} de
voisinage conique. On dit qu'une $X_i\subset X$
est compatible avec un germe, représenté par $\bar T$,
s'il est \inserted{$\exists\lambda\in]0,1[$ \deleted{\unsure{telle}}
tel qu'il soit} compatible avec $\lambda\bar T$
(donc aussi avec $\lambda'\bar T$
avec $0<\lambda'\leq\lambda$) --- ce qui ne dépend pas du
choix du représentant. On a de $=$ la
notion de compatibilité stricte.

Si $\bar T$ et $\bar T'$ sont équivalents, \deleted{et si}
$\bar T$ est
cylindrique ssi $\bar T'$ l'est. On trouve ainsi la notion
d'équivalence de voisinages cylindriques, et celle \inserted{de}
germe de voisinage cylindrique. (Pour qu'il
en existe, il faut s. que $Y$ soit sous-espace bord)
\margin{on dira aussi $Y$ borde $X$, ou dire est sous-espace
bordant\ldots}
Idem pour équisingularité d'un germe de
voisinage conique.

Soit $(\tilde X^*,\dot T)$, avec $\dot T\subset\unsure{\tilde X|\tilde X^*}$ muni d'un
germe de voisinage cylindrique. Soit $\dot T\longrightarrow Y$
application continue, alors
$\deleted{\unsure{f}}X=\tilde X^*\amalg_{\dot T}Y\supset Y$
est muni d'un germe de voisinage \deleted{\unsure{cy}} conique,
qui est isotrivial ssi $\dot T\longrightarrow Y$ est \unsure{fibrant ?|fibrant}.
\margin{21}
